% généré par scripts/06_textes.py
\dosmeta{Année de classement : 375 (directe : nom consulaire lu). Références : ICUR II 6009 = ILCV 775 (em). Rattachement selon l’EDB : \emph{Coem. s.Zenonis}. Conservation selon l’EDB : \emph{Collegio Teutonico}. Support : Plaque de marbre. 38 × 31 cm (h. × l.). Inscription gravée. Signes décrits par l’EDB : chrisme. Fiche EDB : Maria Cristina Grisanzio, première version 2006. Notice du catalogue p.~\pageref{not:183}.}
\txtbloc{Texte transmis (EDCS, état 2022)}{\lnum{1}Depossio Petroni(a)es(!) VI Kal(endas)\par
\lnum{2}Dec(embres) post consulatum\par
\lnum{3}Gratiani Aug(usti) III et Fl(avi) Equiti\par
\lnum{4}qu(a)e vixit virgo annos p(lus) m(inus)\par
\lnum{5}XVIIII in pace est de regio\par
\lnum{6}ne VIIII a lacu cunic(u)li}
\txtbloc{Texte de l’EDB (EDB21312)}{\lnum{1}depossio Petronies VI kal(endas)\par
\lnum{2}dec(embres) post consulatum\par
\lnum{3}Gratiani Aug(usti) III et Fl(avii) Equiti\par
\lnum{4}que vixit virgo annos p(lus) m(inus)\par
\lnum{5}XVIIII in pace est de regio/ne VIIII a lacu cunic⟨u⟩li ((monogramma Christi))\par
\lnum{6}((pavo)) ((olea))}
